\scanpage{194}

% B01 / text / top, bullet and cancelled lines
$\bullet$\quad Notons que le fait que $G$ soit transitif sur l'ens. des repères de $(S,A_0)$, signifie aussi que
\deleted{a) $G$ est transitif sur $S$ (condition qui ne dép que pour \unsure{tt|tout} $s\in S$, $G_s$ est}

% B02 / display / upper, (17), brace at left
\[
\left\{\begin{aligned}
 &\text{a) }G\text{ est transitif sur }S\quad\text{(condition indépendante du choix de }A_0\text{)}\\
 &\text{b) }\forall s\in S\quad\text{(ou pour \uline{un} }s_0\in S,\text{ c'est kif kif)}\\
 &\quad\text{l'opération de }G_s\text{ sur }\mathrm{Et}_{A_0}(s)=E\text{ est}\\
 &\quad\text{transitif sur l'ens des repères de}\\
 &\quad\text{la structure de contour }(E,A_{0E}).
\end{aligned}\right. \tag*{(17)}
\]

% B03 / text / middle, polygonality and first marginal note
% slanted
\margin{Cel signifie donc que $A_0$ est associé à une structure de 2-polyèdre régulier triangulaire, et que l'application $\mathrm{Et}_{A_0}$ est injective, sauf \unsure{ds|dans} le seul cas de l'octaèdre où \ldots}
Je vais supposer que cette structure de contour,
\uline{dans le cas de $A_0$}, est connexe i.e. polygonale,
de sorte que cette condition implique que

% B04 / display / middle, (18) and second marginal note
\margin{\unsure{R\kern-.1em b|Rb}}\[
\begin{aligned}
 D_E&=\mathrm{Im}(G_s\longrightarrow\unsure{\mathfrak S|\mathfrak G}_E)\quad\text{est égal au groupe}\\
 &\mathrm{Aut}(E,A_{0E})\text{ tout entier, qui est isomorphe}\\
 &\text{au groupe diédral }D_n,\text{ et contient un}\\
 &\text{unique sous-groupe }D_E^+\ \inserted{=\Lambda_E}\ \text{cyclique d'ordre}\\
 &n.
\end{aligned}\tag*{(18)}
\]
% slanted
\margin{On suppose les sommets $s\in S$ d'ordre $n$ pour $A_0$, ou ce qui revient au $\overline{\mathrm m}$, que les $E\in\unsure{\mathcal E|\mathscr E|\Sigma}$ sont de cardinal $n$ $(n\geq3)$}

% B05 / text / lower, dependence on the choices
A priori, \inserted{par sa définition} ce groupe dépend du choix de $s\in S$, mais
sa description comme $\mathrm{Aut}(E,A_{0E})$ ne dépend
pas \deleted{\unsure{des|du}} choix de $s$ -- elle dépend par contre de
celui de $A_0$. En fait, on voit qu'il n'en
dépend pas, mais seulement des données

% B06 / display / lower, (19) and note
% slanted
\margin{On pourrait faire jouer un rôle \unsure{symétrique à|symétriquement en} $S$ et $\mathcal E$}
\[
 (S,\ \deleted{\unsure{\text{groupe}}}\,G_0\subset\mathfrak S_S,\ \mathcal E\subset\unsure{P|\mathcal P|\mathfrak P}(S)) \tag*{(19)}
\]

% B07 / text / lower, before (20)
Car on a en effet

% B08 / display / lower, (20)
\[
 D_E=\mathrm{Im}(G_E\longrightarrow\mathfrak S_E), \tag*{(20)}
\]

% B09 / text / foot, continuing
(alors que a priori l'image de $G_E\supset G_s$ dans
$\mathfrak S_E$ contient celle de $G_s$ et \deleted{\unsure{pourrait}}\inserted{semblerait}

% archive sheet number: 193
